\subsection{凸函数族}

命题 2. 设 \(f_{i}\) （ \(1 \leqslant i \leqslant p\) ）是 \(p\) 个凸函数（定义在区间 \(\mathbf{I} \subset \mathbf{R}\) 上的），\(c_{i}\) （ \(1 \leqslant i \leqslant p\) ）是 \(p\) 个任意正数；则函数 \(f = \sum_{i=1}^{p} c_{i} f_{i}\) 在 \(\mathbf{I}\) 上凸。此外，若对至少一个指标 \(j\) ，函数 \(f_{j}\) 在 \(\mathbf{I}\) 上严格凸，且 \(c_{j} > 0\) ，则 \(f\) 在 \(\mathbf{I}\) 上严格凸。

把不等式 (1)（相应地，(2)）应用于每个 \(f_{i}\) ，把关于 \(f_{i}\) 的不等式乘以 \(c_{i}\) ，然后逐项相加，便立即得出结论。

命题 3. 设 \((f_{\alpha})\) 是区间 \(\mathbf{I} \subset \mathbf{R}\) 上的一族凸函数；若此族的上包络 \(g\) 在 \(\mathbf{I}\) 的每一点处有限，则 \(g\) 在 \(\mathbf{I}\) 上凸。

事实上，位于 g 的图像上方的点 \((x, y) \in \mathbf{R}^{2}\) 所成的集是由位于诸函数 \(f_{\alpha}\) 中每一个的图像上方的点所成的凸集之交；因此它是凸的。

命题 4. 设 H 是区间 \(I \subset R\) 上的一集凸函数；若 F 是 H 上的一个滤子，它在 I 上逐点收敛于一个有限实函数 \(f_{0}\) ，则此函数在 I 上凸。

为看出这一点，只需在不等式 (1) 中沿 F 取极限。
% semantic-fragments: legacy-input-seam
\relax{}
